\documentclass[UTF8,fontset=windows,12pt]{ctexart}
%\setcounter{tocdepth}{1}
%\setcounter{secnumdepth}{4}
%\usepackage{ctex} %若要输入中文请取消注释或者把article改成ctexart
% This first part of the file is called the PREAMBLE. It includes
% customizations and command definitions. The preamble is everything
% between \documentclass and \begin{document}.
%\usepackage[OT1]{fontenc}
%\usepackage{fontspec}
%\setmainfont{Times New Roman}
%\setCJKmainfont{Microsoft YaHei}
\usepackage[left=1.91cm,right=1.91cm,top=2.54cm,bottom=2.54cm]{geometry}
\usepackage{amsfonts}
\usepackage{amsmath}
\usepackage{amssymb}
\usepackage[upint,lcgreekalpha]{stix}
\usepackage[type1]{garamondlibre}
%\usepackage{esint}
\usepackage{pgfplots}
\pgfplotsset{compat=1.18}
\usepackage{tikz}
\usetikzlibrary{arrows}
\usepackage{extarrows}
\usepackage[only,llbracket,rrbracket]{stmaryrd}
\usepackage{titlesec}
\usepackage{bm}
\usepackage{graphicx}
\usepackage{appendix}
\usepackage{tikz-cd}
\usepackage{tikz-3dplot}
\usetikzlibrary{arrows.meta, decorations.pathreplacing, calc, patterns}
\usepackage{float}
\usepackage{mathtools}
\usepackage{amsthm}
\usepackage{verbatim}              % better theorem environments
\usepackage{setspace}
\usepackage{enumitem}
\setenumerate[1]{itemsep=0pt,partopsep=0pt,parsep=\parskip,topsep=5pt}
\setitemize[1]{itemsep=0pt,partopsep=0pt,parsep=\parskip,topsep=5pt}
\setdescription{itemsep=0pt,partopsep=0pt,parsep=\parskip,topsep=5pt}
\usepackage{xcolor}
\usepackage[colorlinks,linkcolor=black,anchorcolor=black,citecolor=black]{hyperref}
%\usepackage{apacite}
\usepackage[numbers]{natbib}
%\usepackage{showkeys}
\usepackage{cases}
\usepackage{fancyhdr}
\usepackage[scaled=0.92]{helvet}	% set Helvetica as the sans-serif font
%\usepackage{upgreek}
%\renewcommand{\rmdefault}{ptm}		% set Times as the default text font
%\usepackage{newtxtext}
\bibliographystyle{plain}
% If AMS-LaTeX is used, it can be loaded before or after mtpro2		
% The following loads metro and defines some common MTPro options [2, 4]
%\usepackage[lite,subscriptcorrection,nofontinfo]{mtpro2}
%\pdfmapfile{+mtpro2.map}

%\usepackage{txfonts}

% various theorems, numbered by section

\makeatletter
\renewenvironment{proof}[1][\proofname]{\par
  \pushQED{\qed}%
  \normalfont \topsep6\p@\@plus6\p@\relax
  \trivlist
  \item[\hskip\labelsep
        \bfseries % 关键修改：斜体改为加粗
        #1\@addpunct{.}]\ignorespaces
}{%
  \popQED\endtrivlist\@endpefalse
}
\makeatother
\newtheoremstyle{kaiti-theorem}   % 样式名称
  {3pt}                           % 上方间距
  {3pt}                           % 下方间距
  {\kaishu\upshape}               % 正文字体：中文楷体 + 英文直立
  {}                              % 缩进量
  {\bfseries\upshape}      % 头部字体（定理名）：中文楷体粗体 + 英文直立粗体
  {.}                             % 定理名后标点
  {0.5em}                         % 定理名后间距
  {}                              % 头部说明（留空）

% 应用自定义样式
\theoremstyle{kaiti-theorem}
\newtheorem{thm}{定理}[section]
\newtheorem{nota}[thm]{记号}
\newtheorem{question}{问题}
\newtheorem*{questionn}{思考}
\newtheorem{exe}{习题}[section]
\newtheorem{assump}[thm]{假设}
\newtheorem*{thmm}{定理}
\newtheorem{defn}{定义}[section]
\newtheorem{lem}[thm]{引理}
\newtheorem*{lemm}{引理}
\newtheorem{prop}[thm]{命题}
\newtheorem*{propp}{命题}
\newtheorem*{claim}{断言}
\newtheorem{cor}[thm]{推论}
\newtheorem{conj}[thm]{猜想}
\newtheorem{rmk}{注记}[section]
\newtheorem{ex}{例}[section]
\newtheorem{pf}{证明}
\newtheorem{problem}{作业题}
\numberwithin{equation}{section}





\DeclareMathOperator{\id}{id}
\DeclareMathOperator*{\esssup}{ess\,sup}
\renewcommand{\Re}{\textbf{Re}\,}  
\renewcommand{\Im}{\textbf{Im}\,}  
\newcommand{\R}{\mathbb{R}}  
\newcommand{\FH}{\mathbb{H}}  
\newcommand{\C}{\mathbb{C}}  
\newcommand{\Z}{\mathbb{Z}}
\newcommand{\X}{\mathbb{X}}
\newcommand{\N}{\mathbb{N}}
\newcommand{\Q}{\mathbb{Q}}
\newcommand{\T}{\mathbb{T}}
\newcommand{\TT}{\mathcal{T}}
\newcommand{\TP}{\overline{\partial}{}}
\newcommand{\TJ}{\langle\TP\rangle}
\newcommand{\TL}{\overline{\Delta}{}}
\newcommand{\curl}{\text{curl }}
\newcommand{\dive}{\text{div }}
\newcommand{\bd}[1]{\mathbf{#1}}  % for bolding symbols
\newcommand{\bds}[1]{\boldsymbol{#1}}  % for bolding symbols
\newcommand{\RR}{\mathcal{R}}      % for Real numbers
\newcommand{\sh}{\mathcal{S}}  
\newcommand{\sph}{\mathbb{S}}  
\newcommand{\Om}{\Omega}
\newcommand{\OmT}{{\Omega_T}}
\newcommand{\ut}{{U_T}}
\newcommand{\vp}{\varphi}
\newcommand{\Th}{\Theta}
\newcommand{\Lam}{\Lambda}
\newcommand{\Omb}{{\overline{\Omega}}}
\newcommand{\OmbT}{{\overline{\Omega_T}}}
\newcommand{\ub}{{\overline{U}}}
\newcommand{\ubt}{{\overline{U_T}}}
\newcommand{\q}{\quad}
\newcommand{\p}{\partial}
\newcommand{\pb}{\boldsymbol{\partial}}
\newcommand{\lap}{\Delta}
\newcommand{\dd}{\mathfrak{D}}
\newcommand{\DD}{\mathcal{D}}
\newcommand{\den}{{\bar\rho}_{0}}
\newcommand{\Dt}{{\p}_{t}+v^{k}{\p}_{k}}
\newcommand{\col}[1]{\left[\begin{matrix} #1 \end{matrix} \right]}
\newcommand{\noi}{\noindent}
\newcommand{\nab}{\nabla}
\newcommand{\cnab}{\overline{\nab}}
\newcommand{\cp}{\overline{\partial}{}}
\newcommand{\dx}{\,\mathrm{d}x}
\newcommand{\dxx}{\,\mathrm{d}\bds{x}}
\newcommand{\xxi}{\boldsymbol{\xi}}
\newcommand{\eeta}{\boldsymbol{\eta}}
\newcommand{\dxxi}{\,\mathrm{d}\boldsymbol{\xi}}
\newcommand{\deeta}{\,\mathrm{d}\boldsymbol{\eta}}
\newcommand{\dxi}{\,\mathrm{d}\xi}
\newcommand{\deta}{\,\mathrm{d}\eta}
\newcommand{\dy}{\,\mathrm{d}y}
\newcommand{\dyy}{\,\mathrm{d}\bds{y}}
\newcommand{\dz}{\,\mathrm{d}z}
\newcommand{\dzz}{\,\mathrm{d}\bds{z}}
\newcommand{\dph}{\,\mathrm{d}\varphi}
\newcommand{\dt}{\,\mathrm{d}t}
\newcommand{\dr}{\,\mathrm{d}r}
\newcommand{\dtt}{\,\mathrm{d}\tau}
\newcommand{\dS}{\,\mathrm{d}S}
\newcommand{\dss}{\,\mathrm{d}\sigma}
\newcommand{\dmu}{\,\mathrm{d}\mu}
\newcommand{\dnu}{\,\mathrm{d}\nu}
\newcommand{\dV}{\,\mathrm{d}V}
\newcommand{\ds}{\,\mathrm{d}s}
\newcommand{\drr}{\,\mathrm{d}\rho}
\newcommand{\dist}{\text{dist }}
\newcommand{\vol}{\text{vol}\,}
\newcommand{\volo}{\text{vol}(\Omega)}
\newcommand{\spt}{\text{Spt}\,}
\newcommand{\Tr}{\text{Tr}\,}
\newcommand{\lee}{\langle}
\newcommand{\ree}{\rangle}
\newcommand{\xee}{\langle\xi\rangle}
\newcommand{\xxee}{\langle\boldsymbol{\xi}\rangle}
\newcommand{\xeta}{\langle\eta\rangle}
\newcommand{\xeeta}{\langle\boldsymbol{\eta}\rangle}
\newcommand{\kk}{\kappa}
\newcommand{\lam}{\lambda}
\newcommand{\io}{\int_{\Omega}}
\newcommand{\iu}{\int_{U}}
\newcommand{\is}{\int_{\Sigma}}
\newcommand{\ipo}{\int_{\partial\Omega}}
\newcommand{\ipu}{\int_{\partial U}}
\newcommand{\ig}{\int_{\Gamma}}
\newcommand{\ir}{\int_{\R}}
\newcommand{\ird}{\int_{\R^d}}
\newcommand{\irn}{\int_{\R^n}}
\newcommand{\fpi}{\frac{1}{(2\pi)^{\frac{d}{2}}}}
\newcommand{\ddt}{\frac{\mathrm{d}}{\mathrm{d}t}}
\newcommand{\ddx}{\frac{\mathrm{d}}{\mathrm{d}x}}
\newcommand{\eps}{\varepsilon}
\providecommand{\jump}[1]{\left\llbracket #1 \right\rrbracket }
\providecommand{\len}[1]{\lee #1 \ree }
\providecommand{\ino}[1]{\left\| #1 \right\| }
\providecommand{\bno}[1]{\left| #1 \right| }
\newcommand{\red}{\textcolor{red}}
\newcommand{\green}{\textcolor{green}}
\newcommand{\blue}{\textcolor{blue}}
\newcommand{\nnr}{{[n+1]}}
\newcommand{\nnn}{{[n]}}
\newcommand{\nnl}{{[n-1]}}
\newcommand{\nnll}{{[n-2]}}	
\newcommand{\mmr}{{[m+1]}}
\newcommand{\mmm}{{[m]}}
\newcommand{\mml}{{[m-1]}}
\newcommand{\mmll}{{[m-2]}}

%---------------Special variables--------------
\newcommand{\CC}{\mathfrak{C}}  
\newcommand{\NN}{\mathbf{N}}
\newcommand{\LH}{\mathcal{L}}
\newcommand{\HH}{\mathcal{H}}  
\newcommand{\EE}{\mathcal{E}}  
\newcommand{\FT}{\mathcal{F}}  
\newcommand{\fT}{\mathbf{T}}  
\newcommand{\FA}{\mathbf{A}}  
\newcommand{\MH}{\mathcal{M}}
\newcommand{\NH}{\mathcal{N}}
\newcommand{\PH}{\mathcal{P}}
\newcommand{\VH}{\mathcal{V}}
\newcommand{\XX}{\mathfrak{X}}
\newcommand{\xx}{{\bds{x}}}
\newcommand{\bx}{\mathbf{x}}
\newcommand{\by}{\mathbf{y}}
\newcommand{\bp}{\mathbf{p}}
\newcommand{\bq}{\mathbf{q}}
\newcommand{\pp}{{\bds{p}}}
\newcommand{\qq}{{\bds{q}}}
\newcommand{\uu}{\mathbf{u}}
\newcommand{\ww}{\mathbf{w}}
\newcommand{\vv}{\mathbf{v}}
\newcommand{\yy}{{\bds{y}}}
\newcommand{\zz}{{\bds{z}}}
\newcommand{\zo}{\mathbf{0}}
\newcommand{\ee}{\mathbf{e}}
\newcommand{\fa}{\mathbf{a}}
\newcommand{\fb}{\mathbf{b}}
\newcommand{\fc}{\mathbf{c}}
\newcommand{\ff}{\bds{f}}
\newcommand{\fr}{\mathbf{r}}
\newcommand{\fh}{\mathbf{h}}
\newcommand{\fm}{\mathbf{m}}
\newcommand{\fg}{\mathbf{g}}
\newcommand{\FF}{\mathbf{F}}
\newcommand{\GG}{\mathbf{G}}
\newcommand{\BB}{\mathbf{B}}



\begin{document}
\title{\LARGE{\bf  2026年秋季学期~常微分方程~作业一}}
\author{常微分方程的初等计算、解的存在唯一性定理}
\date{截止时间：2026.9.16下课前（电子版截至当天23:59:59）}

\maketitle

\textbf{第5、9题选做，其它必做。禁止抄袭或使用AI直接生成答案PDF，被抓到者此次作业零分。}\\
A组：交流生、高年级同学以及25级学号 $\leqslant$ PB25000295的同学把作业交给张博洋助教。\\
B组：25级学号 $\geqslant$ PB25000304的同学把作业交给周芾助教。

\begin{problem}[习题1.1.4] 
解方程 $x'(t)=x(t)^a$, 其中 $a=1/5, 1, 2$, 并分别作出相应积分曲线族的草图.
\end{problem}

\begin{problem}[习题1.1.5]
令牵引点 $A$ 沿 $x$ 轴正向运动，并以 $A$ 沿 $x$ 轴的位移 $s$ 为参数，$A(s)=(s,0),s\geqslant0.$ 质点 $B$ 从 $ B(0)=(0,b),~(b>0)$ 出发。假设绳长保持不变，即 $\|A(s)-B(s)\|=b$，其中 $B(s)=(x(t),y(t))$；并且 $B$ 的轨迹切向始终沿着由 $B$ 指向 $A$ 的方向，即存在非负函数 $\lambda(s)$ 使得 $(x'(s),y'(s))=\lambda(s)(s-x(s), -y(s))$。求质点 $B$ 的轨迹方程，即解出 $x(s), y(s)$ 的参数表达式。（注意：题干中的 $\lambda(s)$ 是可以跟 $x(s)$ 一起解出来的，最终答案不能含有 $\lambda(s)$。）
\end{problem}

\begin{problem}[习题1.2.4]
求微分方程 $x'(t)\sin 2t=2(x(t)+\cos t)$ 满足 $t\to\pi/2$ 时仍有界的解。
\end{problem}

\begin{problem}[习题1.2.6]
设 \(p, q\) 都是 \(T\)-周期连续函数, 并且 \(\int_0^T p(s)\ds>0. \) 证明: 方程 \(x'+p(t)x=q(t)\) 的任意解与其唯一 \(T\)-周期解之差在 \(t\to+\infty\) 时趋于零. 
\end{problem}


\begin{problem}[习题1.2.8, 选做] 设$E(t)$是取值恒大于零的光滑函数。对任意$t\geqslant 0$，它满足不等式 $$E'(t)\leqslant C(1+t)^{-\sigma}E(t)^p,~E(0)\leqslant\eps^2,$$ 其中 $\sigma>0, p>1, C>0, 0<\eps\ll 1$. 定义 $\tilde{T}_\varepsilon:=\inf\left\{t>0: E(t)=10\varepsilon^2\right\}$，试讨论 $\tilde{T}_\eps$ 的下界对 $\eps$ 的依赖关系，分情况 $\sigma>1, \sigma=1, 0<\sigma<1$ 讨论。
\end{problem}


\begin{problem}[习题1.3.1]
求解方程 $x'(t)=-x(t)^2-\dfrac{1}{4t^2}$. （提示：先找形如 $x_p(t)=C/t$ 的特解）
\end{problem}

\bigskip \hrule \bigskip

考虑如下形式的一阶常微分方程的初值问题
\begin{equation}\label{2-0: IVP}\tag{2.0.1}
 \xx'(t)=\ff(t, \xx(t)), \qquad \xx(t_0)=\xx_0, 
\end{equation}
其中 \(I\subset\R\) 是含有 \(t_0\) 的区间, 未知函数\(\xx:I\to\R^d\), 给定函数 \(\ff:D\to\R^d\) 的定义域\(D\) 是 \(\R\times\R^d\) 中的开集, 并且 \((t_0, \xx_0)\in D\). 

\begin{problem}[习题2.1.2]
对初值问题\eqref{2-0: IVP}，如果 $\ff$ 满足的条件改为：$\ff$ 连续, 且存在 $[t_0,t_0+a]$ 上的可积函数 (即绝对值的积分是收敛的) $k(t)$ 使得
\[
|\ff(t,\xx)|\leqslant k(t)(1+|\xx|),\quad |\ff(t,\xx_1)-\ff(t,\xx_2)|\leqslant k(t)|\xx_1-\xx_2|.
\] 证明: 存在区间 $[t_0,t_0+\delta]$ 使得对应的 Picard 序列在该区间上一致收敛到\eqref{2-0: IVP}的解.
\end{problem}

\begin{problem}[习题2.1.3] 
对初值问题\eqref{2-0: IVP}，设空间维数 $d=1$, 且 $f(t,x)$ 连续，并且关于 $x$ 变量是单调递减的。证明：该初值问题在 $t\geqslant t_0$ 的解是唯一的。并问：在 $t\leqslant t_0$ 上的解是否唯一？证明结论或者构造反例。
\end{problem}

\begin{problem}[习题2.2.2, 选做]
设空间维数 $d=1$, 在区间 $I=[t_0,t_0+h]$ 上 ($h$定义与 Peano 存在性定理相同) 构造 Tonelli 序列如下：任给正整数 $n$, 令 $t_k=t_0+k\dfrac{h}{n}$, 然后我们从 $[t_0,t_1]$ 到 $[t_1,t_2]$, 再从 $[t_1,t_2]$ 到 $[t_2,t_3]$ 最终从 $[t_{n-2},t_{n-1}]$ 到 $[t_{n-1},t_n]$ 分段递归定义逼近解
\[
x_n(t)=\begin{cases}
x_0& t\in[t_0,t_1]\\
x_0+\int_{t_0}^{t-\frac{h}{n}} f(s,x_n(s))\ds&t\in[t_1,t_0+h]
\end{cases}.
\]　请用 Tonelli 序列和 Arzel\`a--Ascoli 引理证明 Peano 存在性定理。
\end{problem}

\end{document}
